\documentclass[12pt]{article}
\usepackage[utf8]{inputenc}
\usepackage[T1]{fontenc}
\usepackage{hyperref}
\title{Petapa Zapotec Language Resources for K-12 Curriculum Development}
\author{Compiled by Claude Sonnet 5.5 (claude-sonnet-5-5) for Andr\'as Kornai}
\date{}
\begin{document}
\maketitle

\section{Introduction}
Petapa Zapotec (ISO 639-3 zpe, Glottocode peta1247) is spoken in Oaxaca, Mexico. It belongs to the Trans-Yautepec group of Central Zapotec. Wikipedia cites about 8,000 speakers (1990 census), while Joshua Project reports 13,000 for the people group. Wikitongues rates it 6b (threatened) and unwritten. Listed varieties include Santa Mar\'ia Petapa, Santo Domingo Petapa and Santiago Lachigir\'i. No school materials were found.

\section{Primary Educational Resources}
None were found in the searches.

\section{Linguistic and Reference Materials}
\begin{itemize}
\item Wikipedia article: \url{https://en.wikipedia.org/wiki/Petapa_Zapotec}
\item Glottolog entry: \url{https://glottolog.org/resource/languoid/id/peta1247}
\item Endangered Languages Project entry and bibliography: \url{https://endangeredlanguages.com/node/38696}
\item Wikitongues language page: \url{https://wikitongues.org/languages/zpe/}
\end{itemize}

\section{Religious and Community Resources}
\begin{itemize}
\item Joshua Project profile (Bible translation reported as started): \url{https://joshuaproject.net/people_groups/14608/MX}
\end{itemize}

\section{Digital and Audio Resources}
\begin{itemize}
\item Global Recordings Network audio Bible teaching: \url{https://globalrecordings.net/en/language/zpe}
\end{itemize}

\section{Curriculum Development Context}
Little published material exists. A curriculum would require new documentation and orthography work with the community. Many speakers reportedly use only their own language, which would favour mother-tongue instruction.
\end{document}